\answer{ex:12:1}{$0.60$ (без корреляции было бы $0.18$), предел $\sqrt{c/(rT)}\approx0.58$: сто нейронов различают лишь «есть/нет».}
\answer{ex:12:2}{(а)~$10^8$ импульсов, $\approx10\,\text{мДж}$: $0.3\,\%$ от $3$~Дж всего мозга. (б)~$3$~Дж, в $\approx300$ раз больше.}
\answer{ex:12:3}{(а)~$\theta_A\in[2,3)$, $\theta_B\in[1,2)$; чужая фигура даёт не больше $2$ и $1$. (б)~$10\cdot96^2\approx9.2\cdot10^4$.}
\answer{ex:12:4}{(а)~$T_d=0.85\,\tau_a=0.34\,\text{с}$. (б)~$T_d\propto\tau_a$ и не зависит от $I$: $\tau_a\approx3.5\,\text{с}$ (моделирование: $3.1\,\text{с}$).}
\answer{ex:12:5}{$p\approx0.17$. Доля падает медленно: $0.90$, $0.88$, $0.86$, $0.84$, $0.82$, $0.79$ при $N=8$, $12$, $24$, $48$, $96$, $192$.}
\answer{ex:12:6}{Все десять прогонов безошибочны при $\tau_{\text{сл}}=20$, $50$ и $200\ms$ (при $30\ms$ один прогон из десяти сорвался); при $5$--$10\ms$ --- около $0.5$, при $0.5$--$2\,\text{с}$ качество падает до $0.86$. Сверху: след должен угаснуть за паузу, $e^{-0.3\,\text{с}/\tau_{\text{сл}}}\ll1$.}
\answer{ex:12:7}{Одной задержки мало: окно $\pm5\ms$ не покрывает $\Delta x/v$ от $5$ до $20\ms$. Две ветви торможения с $5<d_1\le10\ms$ и $15\le d_2\le d_1+10\ms$.}
\answer{ex:12:8}{\texttt{two\_stage}: $\Delta=2$, ничья от одного переворота, смена ответа --- от двух. Счётчик единиц: один переворот; $0.57$ при $p=0.1$.}
